% Part: sets-functions-relations
% Chapter: sets
% Section: basics

\documentclass[../../../include/open-logic-section]{subfiles}

\begin{document}

\olfileid{sfr}{set}{bas}
\olsection{विस्तारात्मकता}

वस्तूंचा समूह एकच वस्तू {म्हणून} विचारात घेतला, की त्याला \emph{संच} म्हणतात.
संच बनवणाऱ्या वस्तूंना त्या संचाचे \emph{घटक} किंवा \emph{सदस्य} म्हणतात.
$x$ हा $A$ या संचाचा !!a{element} असेल, तर आपण $x \in A$ असे लिहितो;
नसेल, तर $x \notin A$ असे लिहितो. ज्या संचात एकही !!{element}s नसतो,
त्याला \emph{रिक्त} संच म्हणतात आणि तो ``$\emptyset$'' या चिन्हाने दर्शवतात.

\begin{explain}
आपण संच \emph{कसा निर्दिष्ट करतो}, त्यातील !!{element}s \emph{कोणत्या क्रमाने}
मांडतो किंवा तेच !!{element}s \emph{किती वेळा} मोजतो, याला महत्त्व नसते.
त्यात कोणते !!{element}s आहेत, एवढेच महत्त्वाचे असते. ही गोष्ट आपण पुढील
तत्त्वात नेमकेपणाने मांडतो.
\end{explain}

\begin{defn}[विस्तारात्मकता]
  $A$ आणि $B$ हे संच असतील, तर $A = B$ तेव्हा आणि केवळ तेव्हाच, जेव्हा
  $A$ मधील प्रत्येक !!{element} हा $B$ चाही !!a{element} असतो, आणि उलटही.
\end{defn}

विस्तारात्मकतेमुळे काही चिन्हपद्धती वापरणे योग्य ठरते. सर्वसाधारणपणे,
$a_{1}$, \dots, $a_{n}$ या वस्तू दिलेल्या असतील, तर
$\{a_{1}, \dots, a_{n}\}$ म्हणजे ज्यातील !!{element}s हे
$a_1, \ldots, a_n$ आहेत \emph{तो विशिष्ट} संच होय. येथे ``\emph{तो विशिष्ट}''
या शब्दांवर भर आहे, कारण विस्तारात्मकतेनुसार असा \emph{एकच} संच असू शकतो.
विस्तारात्मकतेमुळे पुढील समानताही योग्य ठरते:
  \[
    \{a, a, b\} = \{a, b\} = \{b,a\}.
  \]
म्हणजेच, संचांचा विचार करताना त्यातील !!{element}s कोणत्या क्रमाने किंवा
किती वेळा निर्दिष्ट केले आहेत, याला महत्त्व नसते.

\begin{tagblock}{novice}
\begin{ex}
काही वस्तू दिलेल्या असतील, तर त्या एकत्र करून त्यांचा संच बनवता येतो.
उदाहरणार्थ, रिचर्डच्या भावंडांच्या संचात एक व्यक्ती आहे, आणि तो
$S=\{\textrm{Ruth}\}$ असा लिहिता येईल.
$4$ पेक्षा लहान धन पूर्णांकांचा संच $\{1, 2, 3\}$ आहे; पण तो
$\{3, 2, 1\}$ किंवा अगदी $\{1, 2, 1, 2, 3\}$ असाही लिहिता येतो.
विस्तारात्मकतेनुसार हे सर्व एकच संच आहेत. कारण $\{1, 2, 3\}$ मधील
प्रत्येक !!{element} हा $\{3, 2, 1\}$ चा (आणि $\{1,
2, 1, 2, 3\}$ चाही) !!a{element} आहे, आणि उलटही.
\end{ex}
\end{tagblock}

अनेकदा संचातील सर्व !!{element}s यांना समान असलेला एखादा गुणधर्म सांगून आपण
संच निर्दिष्ट करू. त्यासाठी पुढील संक्षिप्त चिन्हपद्धती वापरू:
$\Setabs{x}{\phi(x)}$. येथे $\phi(x)$ हा असा गुणधर्म दर्शवतो की,
$x$ ची संचातील !!{element}s मध्ये गणना होण्यासाठी त्यात तो गुणधर्म असला पाहिजे.

\begin{tagblock}{novice}
\begin{ex}
आपल्या उदाहरणातील $S$ हा संच पुढीलप्रमाणेही निर्दिष्ट करता आला असता:
\[
S = \Setabs{x}{x \text{ हे रिचर्डचे भावंड आहे}}.
\]
\end{ex}
\end{tagblock}

\begin{tagblock}{math}
\begin{ex}
एखादी संख्या तिच्या स्वतःखेरीज इतर सर्व धन विभाजकांच्या बेरजेइतकी असेल तेव्हा
आणि केवळ तेव्हाच तिला \emph{परिपूर्ण} म्हणतात (हे विभाजक म्हणजे त्या संख्येला
निःशेष भाग जाणाऱ्या, पण तिच्याशी समान नसलेल्या संख्या).
उदाहरणार्थ, $6$ ही परिपूर्ण संख्या आहे, कारण तिचे असे विभाजक
$1$, $2$ आणि~$3$ आहेत, आणि $6 = 1 + 2 + 3$.
खरे तर, $10$ पेक्षा लहान धन पूर्णांकांपैकी $6$ हीच एकमेव परिपूर्ण संख्या आहे.
म्हणून विस्तारात्मकतेचा उपयोग करून आपण असे म्हणू शकतो:
\[
	\{6\} = \Setabs{x}{x\text{ ही परिपूर्ण संख्या आहे आणि }0 \leq x \leq 10}
\]
उजवीकडील चिन्हलेखनाचे वाचन असे करतो: ``$x$ ही परिपूर्ण संख्या आहे आणि
$0 \leq x \leq 10$ अशा सर्व $x$ चा संच.'' संच कसा निर्दिष्ट केला आहे,
याला संचाचा विचार करताना महत्त्व नसते, ही गोष्ट वरील समानता स्पष्ट करते.
अधिक सर्वसाधारणपणे, $\phi(x)$ असणाऱ्या सर्व $x$ चा संच नेहमी एकच असतो,
याची विस्तारात्मकता खात्री देते.
त्यामुळे $\Setabs{x}{\phi(x)}$ याला $\phi(x)$ असणाऱ्या सर्व $x$ चा
\emph{तो विशिष्ट} संच म्हणणे विस्तारात्मकतेमुळे योग्य ठरते.
\end{ex}
\end{tagblock}

संच समान असल्याचे दाखवण्याची पद्धत विस्तारात्मकतेमुळे मिळते:
$A = B$ हे दाखवण्यासाठी, $x \in A$ असेल तेव्हा $x \in B$ देखील असते,
आणि $y \in B$ असेल तेव्हा $y \in A$ देखील असते, हे दाखवा.

\begin{prob}
जास्तीत जास्त एकच रिक्त संच असतो हे सिद्ध करा. म्हणजे, $A$ आणि $B$ या
संचांत एकही !!{element}s नसेल, तर $A = B$ हे दाखवा.
\end{prob}

\end{document}
